% ============================================================
\section{10.3 분류를 넘어서: 탐지, 분할, 전이학습}
\begin{frame}{분류를 넘어서: 두 가지 새 문제}
  지금까지의 CNN은 ``이미지 \textbf{전체}가 어떤 클래스인가''(분류)만
  답했다. 실전 문제는 더 세밀한 답을 요구한다:
  \begin{columns}[T]
    \begin{column}{0.5\textwidth}
      \kb{객체 탐지(Object Detection)}
      \begin{itemize}
        \item ``이미지 안에 사물이 \textbf{어디에}, \textbf{무엇}이
          있는가''
        \item 각 사물의 \textbf{위치}를 사각형
          (bounding box, 4개 좌표)으로, 그 안의 \textbf{클래스}를
          함께 예측
      \end{itemize}
    \end{column}
    \begin{column}{0.5\textwidth}
      \kb{의미론적 분할(Semantic Segmentation)}
      \begin{itemize}
        \item ``\textbf{각 픽셀}이 어떤 클래스에 속하는가''
        \item 입력과 \textbf{동일한 크기}의 출력을 내며 각 픽셀에
          클래스 라벨을 붙임
        \item 자율주행: ``이 픽셀은 도로/보행자/차량''
      \end{itemize}
    \end{column}
  \end{columns}
  \vskip0.6em
  \kq{탐지는 ``\textbf{위치까지 붙은 분류}''다 ---
  분류는 벡터 1개를 내지만, 탐지는 사물 \textbf{마다}
  (클래스, x, y, w, h)를 내야 하고 \textbf{사물의 개수조차 모름}.}
  어려운 부분은 ``무엇을 보는가''(10.2가 이미 잘한다)가 아니라
  ``여러 개, 각기 다른 위치''를 \textbf{한 번의 전진전파}로
  다루는가.
\end{frame}
